The 2WikiMultiHopQA-ONCU-500 results provide an additional realistic multi-hop validation setting. Across all three models, full-context input outperforms retrieved-evidence input on relaxed answer F1. Qwen2.5-14B obtains 0.549 relaxed F1 under full context compared with 0.378 under retrieved evidence; Qwen3-14B obtains 0.560 compared with 0.400; and Gemma3-12B obtains 0.537 compared with 0.395. Evidence F1 shows the same direction: full-context evidence F1 ranges from 0.594 to 0.663, whereas retrieved-evidence evidence F1 ranges from 0.316 to 0.352. The ONCU-Relaxed-F1 results preserve this pattern. Full-context ONCU is 0.610 for Qwen2.5-14B, 0.534 for Qwen3-14B, and 0.369 for Gemma3-12B, while retrieved-evidence ONCU is 0.381, 0.367, and 0.215, respectively. Failure-type analysis explains the gap: retrieved evidence produces substantially more evidence-localization failures, ranging from 40.6\% to 44.0\%, whereas full-context localization failures remain between 1.2\% and 10.4\%. These results support the interpretation that, in realistic multi-hop settings, deterministic lexical retrieve-then-read evaluation can become bottlenecked by evidence-chain coverage, even when the full-context condition still imposes a non-trivial integration burden on the model.